\scanpage{38}
\uline{Automorphism of order two of a holomorphic sphere} \inserted{$\Sigma$}
\uline{\unsure{with|:}} \inserted{a $K$ of} $5$ \uline{marked points} : Exactly one fixed point in $K$,
let $K=\{P\}\amalg I$, and $\Sigma$ is equipped with $I$ and an automorphism
\inserted{$\sigma$} of order $2$, stabilizing $I$,
\inserted{$I$ containing no fixed point of $\sigma$.}
Thus this amounts to the datum

a) Of a holomorphic sphere $\Sigma$, an automorphism \inserted{$\sigma$} of order $2$
and a marked fixed point --- calling this one $\infty$, and the other fixed point $0$,
this amounts to the datum of a torsor $T_1\ (= \Sigma\setminus\{0,\infty\})$
under $\mathbb C^*$, i.e. of a complex vector line $L_1$
$(T_1=L_1^*)$

b) Of a pair of \inserted{(distinct)} orbits of
$\{\pm1\}\subset\mathbb C^*$ acting on $T_1=L_1^*$.
Introducing $T_2=T^{\wedge2}$\inserted{$\simeq L_2^*$, where
$L_2=L_1^{\otimes2}$,}, this amounts to the datum of a pair of points
\inserted{$\{\xi,\eta\}$} (distinct) of $T_2$.
Consider then the point
\[
 e=\xi\wedge\eta=\eta\wedge\xi
 \in T_4=T^{\wedge4}\simeq L_4^*
 \quad(\text{where }L_4=L_1^{\otimes4})
\]
we therefore have an
\[
 \mathrm{iso}\quad T^{\wedge4}\simeq\mathbf1_{\mathbb C^*}
\]
whence a reduction of the structure group \inserted{of $T$}
from $\mathbb C^*$ : $\mu_4(\mathbb C_*)\simeq\mathbb Z/4\mathbb Z$.
One can say that the datum of $\{\xi,\eta\}$ amounts to that of
$e=\xi\wedge\eta\inserted{\in T_4}$ and of
$\xi+\eta\in\unsure{t_2|T_2}$, subject to the condition that
$(\xi\unsure{\wedge|+}\eta)^2-4\xi\wedge\eta\ne0$ \deleted{\unsure{a}}.
Note \deleted{$\unsure{\xi^2+\eta^2-2\xi\wedge\eta\ne}$}
that $T_2$ admits a reduction of the structure group $\mathbb C^*$ :
$\{\pm1\}$, by the bibasis $\pm\sqrt e$, thus $\xi+\eta$
defines an \unsure{element} of $\mathbb C$ up to sign, whose
\deleted{$\unsure{(\xi+\eta)^2}$}
\margin{37}
